% ======================================================================
% APPENDIX (v3): the crisis and discovery checks, derived from the ledger.
% \input from changing_networked_mind_v3.tex (after the operator appendix).
% Uses the main paper's macros (\Top, \KL, \Eq) and labels (eq:tilted,
% eq:bmr, eq:gamma, sec:structurelearning, sec:revision, tab:dials).
% ======================================================================

% ----------------------------------------------------------------------
\section{The crisis and discovery checks, derived from the ledger}
\label{app:crisis-discovery}
% ----------------------------------------------------------------------

The per-round loop of \S\ref{sec:model} contains two threshold events: the \emph{crisis check}
(when does an agent prune the protective edges and release the hub?) and the \emph{discovery
check} (when does it wake the unconceived node?). This appendix states exactly what each check
computes, derives both from the single objective the paper already carries, and lists the
modelling commitments that remain after the derivation. The aim is that nothing in the two
checks should read as a stipulated event: every acceptance is the model's own evidence weighed
on the model's own ledger, and everything that is \emph{not} derived is named as a commitment.

\subsection{One ledger for every structural move}

The bounded agent of \S\ref{sec:revision} revises its beliefs by minimizing
$\KL{q}{p(x\mid o,G)}-\lambda\,\Eq{q}{U^{\!\top}\!x}$ (Eq.~\ref{eq:tilted}): fidelity to the
honest posterior, traded against the value of where the commitments land, at temperature
$\lambda$. A structural move---pruning an edge, waking a node---changes both terms, so the same
objective scores the move:
\begin{equation}
\mathrm{score}(M\!\to\!\tilde M)
\;=\;
\underbrace{\Delta F}_{\text{log Bayes factor}}
\;+\;
\underbrace{\Delta G}_{\text{epistemic gain}}
\;+\;
\lambda\,\underbrace{\Delta U}_{\text{conviction value}},
\qquad
\text{accept}\iff \mathrm{score}>0,
\label{eq:ledger}
\end{equation}
with $\Delta F$ the evidence for the edited structure, $\Delta G$ the expected information gain
about any newly introduced latent (zero for a prune), and
$\Delta U = \Eq{q_{\tilde M}}{U^{\!\top}\!x}-\Eq{q_{M}}{U^{\!\top}\!x}$ the change in the
conviction-field value of the commitments. Equation~\eqref{eq:ledger} is not a new mechanism:
it is Eq.~\eqref{eq:tilted} evaluated at two structures and differenced. Both threshold events
of the simulation are instances of it.

\subsection{The crisis check is the ledger on a prune}

For a prune, every term of Eq.~\eqref{eq:ledger} is closed-form in the linear-Gaussian setting.
The evidence half is Bayesian model reduction (Eq.~\ref{eq:bmr}): writing the agent's
accumulated likelihood deposit as $J=\Pi-\Pi_0$, $j=h-h_0$, the Savage--Dickey identity scores
the swap of the prior $(\Pi_0,h_0)$ for the edge-$e$-reduced prior $(\Pi_0^{(e)},h_0^{(e)})$
under the same likelihood, and---this is the part the main text does not spell out---it
\emph{already exhibits the reduced posterior},
\begin{equation}
\Pi^{(e)} = \Pi_0^{(e)} + J,
\qquad
h^{(e)} = h_0^{(e)} + j ,
\end{equation}
so the value half costs one linear solve per candidate edge:
\begin{equation}
\Delta U_e \;=\; U^{\!\top}\!\bigl(\mu^{(e)}-\mu\bigr),
\qquad
\mu^{(e)} = \bigl(\Pi^{(e)}\bigr)^{-1}h^{(e)},\quad \mu = \Pi^{-1}h,
\label{eq:dU-prune}
\end{equation}
with $U=\Top u$ the agent's \emph{current} conviction field, read along its present couplings.
The crisis check is then exactly the ledger: edge $e$ is flagged iff
$\Delta F_e + \lambda\,\Delta U_e > 0$, and \emph{crisis} is the round at which both protective
belt edges are flagged, upon which the agent adopts the reduced prior composed over every edge
its own ledger flags (the hub released). Pruning a conviction-protective edge moves the
posterior mean away from the cherished commitments, so $\Delta U_e<0$ and the rule reproduces
the intuitive form ``the evidence must overcome $\lambda$ times the value the edge protects.''

A natural surrogate suggests itself here: the static threshold $\Delta F_e > \lambda\,v_e$ with
$v_e = |U_{\mathrm{parent}}|+|U_{\mathrm{child}}|$, the summed conviction magnitude at the
edge's endpoints---available without a solve, and with the
right \emph{form} (it is Eq.~\eqref{eq:ledger} with $\Delta U_e \equiv -v_e$). But it has the
wrong \emph{content}: $v_e$ is a static stand-in for a quantity the model can compute exactly,
so we compute it. The difference is not cosmetic, and its direction is itself informative:
because the reduced and full posteriors share the likelihood deposit, the mean displacement of
a prune---and with it $\lambda\,\Delta U_e$---\emph{shrinks as evidence accumulates}. A
data-rich paradigm is barely protected at the threshold; almost all of its persistence must
come from upstream, from the conviction gate that decides which evidence is allowed to
accumulate at all (Eq.~\ref{eq:gamma}). The static surrogate would quietly overstate threshold
protection; the exact ledger locates lock-in where the paper's results put it: gating plus
isolation, not a high bar at the moment of decision.

\subsection{The discovery check is the ledger on an expansion}

Expansion cannot be closed-form (\S\ref{sec:structurelearning}): a new node needs a real
inversion, and---Stanford's problem of unconceived alternatives---the agent cannot plan toward
a dimension it has not conceived. The discovery check therefore has exactly two parts, and only
two:

\begin{enumerate}
\item \textbf{Arrival.} Candidate proposals arrive at a Poisson rate $r$. This is the one
place a stipulated process enters, and it is a deliberate modelling commitment: a rate is the
maximally agnostic model of conceiving what one has not conceived. The proposal's
\emph{direction} is read from the model's own prediction errors: a hidden hub $b$ wired to the
visible nodes with couplings $c$ leaves, on marginalization, the rank-one footprint
$-cc^{\!\top}/\Pi_{bb}$, so the leading eigenvector of the windowed error second moment is the
best single-hub explanation of the residual. The window is an estimator horizon, not a gate.
\item \textbf{Acceptance.} The arrived proposal is accepted iff the \emph{model log Bayes
factor} of the bordered (hub-wired) model against the agent's current one is positive---that
is, iff the residual genuinely loads on the proposed hub, equivalently iff the immediate
Savage--Dickey prune of the new couplings would be rejected. The epistemic gain $\Delta G$ sits
on the ledger too, but on the other side of the move: being positive for \emph{any} new latent,
it is what pays for attempting the (in general expensive) inversion, and it cannot keep a hub
the data do not hold---an accept test that included it would wake on pure noise. Entertaining
is paid by curiosity; keeping is paid by evidence. One consequence is owned rather than hidden:
repeated looks give the Bayes factor a small false-discovery rate (rare
noise crossings of $\Delta F>0$ in the null world). We regard this as a feature of the honest
model---communities do occasionally adopt spurious structure---and its principled remedy is not
a reinstated threshold but the wake-then-prune cycle named in the limitations, in which a
spuriously woken node is pruned back by the same reduction that scored it.
\end{enumerate}

A tempting third part is deliberately absent: a \emph{trigger}, requiring the residual floor
(the eigenvalue above) to clear some height, perhaps for a sustained number of steps, before
the Bayes factor is consulted. Such a gate would duplicate the accept test's job with a free
constant---the Bayes factor already rejects proposals the residual does not support---and its
only principled reading, a budget on when to pay for the expensive inversion, is a statement
about real scientific communities rather than about this simulation, where the inversion is a
small log-determinant. The discovery check therefore has no tuned threshold: the Poisson rate
carries Stanford's commitment, the eigenvector carries the direction, and the ledger decides.
(The floor is still recorded as telemetry, because it is the visible signature of the residual
the old paradigm had been absorbing.)

\subsection{What remains a modelling commitment}

The derivation above leaves a short, explicit list of choices that are \emph{not} forced by the
formalism. We state them because the paper's claim that ``crisis, revolution and discovery are
threshold crossings of the model's own evidence'' is honest only if its boundary is drawn:

\begin{itemize}
\item \textbf{The belt conjunction.} ``Crisis'' is operationalized as \emph{both} named
mass-law edges flagged in the same round. The per-edge accept is derived; the conjunction over
those two edges is our operational definition of having eliminated the phlogiston mass-law
reading.
\item \textbf{The ordering --- a result, not a commitment.} The simulation attempts expansion
only after an agent's own crisis, but the ordering is enforced by the evidence rather than by
fiat: running the identical proposal + ledger pipeline for every pre-crisis agent at every step
(without applying it), the ledger rejects $13{,}180$ of $13{,}185$ pre-crisis agent-steps; the
five crossings are marginal ($\Delta F < 0.005$, two orders below the post-crisis acceptance
evidence) and are exactly the repeated-look false-positive rate owned above, whose remedy is
an audited wake-then-prune cycle. The intact paradigm absorbs the
residual; there is nothing for the check to find until reduction exposes it.
\item \textbf{The Poisson rate $r$} (swept), the prediction-error \emph{window} (an estimator
horizon), a short \emph{warm-up} before reduction is permitted (an initialization guard), and
the released hub's self-precision are quantitative choices of the same standing as the other
dials of Table~\ref{tab:dials}.
\end{itemize}

% ----------------------------------------------------------------------
